\section{Conclusion}
\label{sec:conclusion}

A fingerprint of relief and texture, measured with one code path on the Moon and on Earth, turns
the analog question into a set of auditable distances. For the Connecting Ridge target it
identifies one candidate window at Holuhraun among \ComparedCount{} evaluated windows, at a relief
distance of \BestReliefDistance{} against a median of \ReliefDistanceMedian{} over the compared
windows. The radar audit documents what a radar comparison would
need: a verified handedness convention on both sides, per pixel incidence and a qualified pair.
The encoder, trained with a physics target and gradient reversal across three bodies, is held to a
promotion rule fixed before training, and until it passes, it remains a labelled second opinion.

The results are candidates, and field work is what would test them. Surveyed topography of the
candidate window at metre and decimetre resolution would show whether the hectometre agreement
extends to the scales that matter for surface operations. Clast size and coverage counts, and
notes on moisture and material, would connect the elevation statistics to the ground. A repeat
radar acquisition under controlled geometry, compared with Mini-RF under a declared convention,
would test whether the scattering signatures agree. Until such observations exist, the shortlist
is a plan for where to look, and the Terraluna application lets anyone measure another Earth
location with the same fingerprint, as a live estimate from public elevation data rather than a
site assessment.
